% Generated by eval/scripts/write_collinearity_explanation.py.
\paragraph*{What changes in the sensitivity fits?} The single-predictor fits remove SPEC from the distance model, or distance from the SPEC model, while retaining the model and language indicators. The pooled comparisons use the same 32,090 scored responses and the same predictor standardization, so their differences reflect the changed adjustment rather than a changed cohort. The distance coefficient changes from $-0.112$ in the joint fit to $-0.385$ without SPEC; the SPEC coefficient changes from $0.954$ to $1.157$ without distance (Table~\ref{tab:graded-single}). Both retain their signs, but their magnitudes change. These are coefficients on the logit of expected score, not direct score increments.

SPEC's VIF remains 7.99 even when distance is omitted, showing that overlap with the retained model and language indicators is still substantial. In the decomposition, IF has VIF 38.41 and CONS has VIF 5.04; separating their coefficients also requires caution. Removing a predictor changes the question being estimated and can introduce omitted-variable bias. These checks therefore describe sensitivity; they do not recover a causal effect or guarantee an isolated predictor contribution.

\paragraph*{Implication for H\textsubscript{3}.} All reported fits pass convergence, finite-estimate and positive-definite covariance checks. Those establish numerical usability, not that the predictors provide independent explanations. Although standardization puts their coefficients on comparable relative scales, a larger absolute coefficient or smaller $p$-value alone cannot establish a stronger, more stable association. Interpretation must also consider predictor overlap, the changes under single-predictor adjustment, and the cohort sensitivities reported below.
